\chapter{2014年上海卷（秋理）}

\section{填空题}

\begin{problem}
函数 \( y = 1 - 2\cos^2(2x) \) 的最小正周期是\fillinblank{}.
\end{problem}

\begin{problem}
若复数 \(z = 1 + 2\i\)，其中 \(\i\) 是虚数单位，则 \(\left(z + \frac{\i}{z}\right) \cdot \overline{z} =\)\fillinblank{}.
\end{problem}

\begin{problem}
若抛物线 \( y^{2}=2px \) 的焦点与椭圆 \( \frac{x^{2}}{9}+\frac{y^{2}}{5}=1 \) 的右焦点重合，则该抛物线的准线方程为\fillinblank{}.
\end{problem}

\begin{problem}
设 \( f(x) = \begin{cases} x, & x \in (-\infty, a), \\ x^2, & x \in [a, +\infty), \end{cases} \) 若 \( f(2) = 4 \)，则 \( a \) 的取值范围为\fillinblank{}.
\end{problem}

\begin{problem}
若实数 \(x\)，\(y\) 满足 \(xy = 1\)，则 \(x^2 + 2y^2\) 的最小值为\fillinblank{}.
\end{problem}

\begin{problem}
若圆锥的侧面积是底面积的 3 倍，则其母线与底面夹角的大小为\fillinblank{}.（结果用反三角函数值表示）
\end{problem}

\begin{problem}
已知曲线 \(C\) 的极坐标方程为 \( \rho(3\cos\theta-4\sin\theta)=1 \)，则 \(C\) 与极轴的交点到极点的距离是\fillinblank{}.
\end{problem}

\begin{problem}
设无穷等比数列 \(\{a_{n}\}\) 的公比为 \(q\)，若 \(a_{1} = \lim\limits_{n\to \infty}(a_{3} + a_{4} + \dots + a_{n})\)，则 \(q =\)\fillinblank{}.
\end{problem}

\begin{problem}
若 \( f(x) = x^{\frac{2}{3}} - x^{-\frac{1}{2}} \)，则满足 \( f(x) < 0 \) 的 \( x \) 的取值范围是\fillinblank{}.
\end{problem}

\begin{problem}
为强化安全意识，某商场拟在未来的连续 10 天中随机选择 3 天进行紧急疏散演练. 则选择的 3 天恰好为连续 3 天的概率是\fillinblank{}.（结果用最简分数表示）
\end{problem}

\begin{problem}
已知互异的复数 \(a\)，\(b\) 满足 \(ab \neq 0\)，集合 \(\{a, b\} = \{a^2, b^2\}\)，则 \(a + b =\)\fillinblank{}.
\end{problem}

\begin{problem}
设常数 \(a\) 使方程 \(\sin x + \sqrt{3}\cos x = a\) 在区间 \([0, 2\pi]\) 上恰有三个解 \(x_1\)，\(x_2\)，\(x_3\)，则 \(x_1 + x_2 + x_3 =\fillinblank{}\).
\end{problem}

\begin{problem}
某游戏的得分为 1， 2， 3， 4， 5，随机变量 \(\xi\) 表示小白玩该游戏的得分. 若 \(E(\xi) = 4.2\)，则小白得 5 分的概率至少为\fillinblank{}.
\end{problem}

\begin{problem}
已知曲线 \(C: x = -\sqrt{4 - y^2}\)，直线 \(l: x = 6\). 若对于点 \(A(m,0)\)，存在 \(C\) 上的点 \(P\) 和 \(l\) 上的 \(Q\) 使得 \(\overrightarrow{AP} + \overrightarrow{AQ} = \overrightarrow{0}\)，则 \(m\) 的取值范围为\fillinblank{}.
\end{problem}

\section{单选题}

\begin{problem}
设 \(a\)，\(b \in \R\)，则“\(a + b > 4\)”是“\(a > 2\) 且 \(b > 2\)”的（\quad）

\choices
    {充分非必要条件}
    {必要非充分条件}
    {充要条件}
    {既不充分也不必要条件}
\end{problem}

\begin{problem}
如图，四个棱长为 1 的正方体排成一个正四棱柱，\(AB\) 是一条侧棱， \( P_{i} \) \( (i=1,2,\cdots,8) \) 是上底面上其余的八个点，则 \( \overrightarrow{AB} \cdot \overrightarrow{AP_{i}} \) \( (i=1,2,\cdots,8) \) 的不同值的个数为（\quad）

\bitmapfigure[max width=0.42\linewidth]{img/2014/shanghai_science/q16_fig1.png}

\choices
    {1}
    {2}
    {4}
    {8}
\end{problem}

\begin{problem}
已知 \(P_{1}(a_{1}, b_{1})\) 与 \(P_{2}(a_{2}, b_{2})\) 是直线 \(y = kx + 1\)（\(k\) 为常数）上两个不同的点，则关于 \(x\) 和 \(y\) 的方程组 \(\begin{cases} a_{1}x + b_{1}y = 1, \\ a_{2}x + b_{2}y = 1 \end{cases}\) 的解的情况是（\quad）

\choices
    {无论 \(k, P_{1}, P_{2}\) 如何，总是无解}
    {无论 \(k\)，\(P_{1}\)，\(P_{2}\) 如何，总有唯一解}
    {存在 \(k\)，\(P_{1}\)，\(P_{2}\)，使之恰有两解}
    {存在 \(k\)，\(P_{1}\)，\(P_{2}\)，使之有无穷多解}
\end{problem}

\begin{problem}
设 \( f(x) = \begin{cases} (x - a)^2, & x \leqslant 0, \\ x + \frac{1}{x} + a, & x > 0, \end{cases} \) 若 \( f(0) \) 是 \( f(x) \) 的最小值，则 \( a \) 的取值范围为（\quad）

\choices
    {\([-1, 2]\)}
    {\([-1,0]\)}
    {\([1, 2]\)}
    {\([0,2]\)}
\end{problem}

\section{解答题}

\begin{problem}
底面边长为 2 的正三棱锥 \(P-ABC\)，其表面展开图是三角形 \(P_{1}P_{2}P_{3}\)，如图，求 \( \triangle P_{1}P_{2}P_{3} \) 的各边长及此三棱锥的体积 \(V\).

\bitmapfigure[max width=0.34\linewidth]{img/2014/shanghai_science/q19_fig1.png}
\end{problem}

\begin{problem}
设常数 \(a \geqslant 0\)，函数 \(f(x) = \frac{2^x + a}{2^x - a}\).

\begin{enumerate}
\item 若 \(a = 4\)，求函数 \(y = f(x)\) 的反函数 \(y = f^{-1}(x)\)；
\item 根据 \(a\) 的不同取值，讨论函数 \(y = f(x)\) 的奇偶性，并说明理由.
\end{enumerate}
\end{problem}

\begin{problem}
如图，某公司要在 \(A\)，\(B\) 两地连线上的定点 \(C\) 处建造广告牌 \(CD\)，其中 \(D\) 为顶端，\(AC\) 长 35 米，\(CB\) 长 80 米. 设点 \(A\)，\(B\) 在同一水平面上，从 \(A\) 和 \(B\) 看 \(D\) 的仰角分别为 \(\alpha\) 和 \(\beta\).

\begin{enumerate}
\item 设计中 \(CD\) 是铅垂方向. 若要求 \(\alpha \geqslant 2\beta\)，问 \(CD\) 的长至多为多少（结果精确到 0.01 米）？
\item 施工完成后，\(CD\) 与铅垂方向有偏差. 现在实测得 \(\alpha = 38.12^{\circ}\)，\(\beta = 18.45^{\circ}\)，求 \(CD\) 的长（结果精确到 0.01 米）.
\end{enumerate}

\bitmapfigure[max width=0.56\linewidth]{img/2014/shanghai_science/q21_fig1.png}
\end{problem}

\begin{problem}
在平面直角坐标系 \(xOy\) 中，对于直线 \(l: ax + by + c = 0\) 和点 \(P_{1}(x_{1}, y_{1})\)，\(P_{2}(x_{2}, y_{2})\)，记 \(\eta = (ax_{1} + by_{1} + c)(ax_{2} + by_{2} + c)\). 若 \(\eta < 0\)，则称点 \(P_{1}\)，\(P_{2}\) 被直线 \(l\) 分隔. 若曲线 \(C\) 与直线 \(l\) 没有公共点，且曲线 \(C\) 上存在点 \(P_{1}\)，\(P_{2}\) 被直线 \(l\) 分隔，则称直线 \(l\) 为曲线 \(C\) 的一条分隔线.

\begin{enumerate}
\item 求证：点 \(A(1,2)\)，\(B(-1,0)\) 被直线 \(x + y - 1 = 0\) 分隔；
\item 若直线 \(y = kx\) 是曲线 \(x^2 - 4y^2 = 1\) 的分隔线，求实数 \(k\) 的取值范围；
\item 动点 \(M\) 到点 \(Q(0,2)\) 的距离与到 \(y\) 轴的距离之积为 1，设点 \(M\) 的轨迹为 \(E\)，求证：通过原点的直线中，有且仅有一条直线是 \(E\) 的分割线.
\end{enumerate}
\end{problem}

\begin{problem}
已知数列 \(\{a_{n}\}\) 满足 \(\frac{1}{3} a_{n} \leqslant a_{n+1} \leqslant 3a_{n}\)，\(n \in \mathbf{N}^{*}\)，\(a_{1} = 1\).

\begin{enumerate}
\item 若 \(a_{2} = 2\)，\(a_{3} = x\)，\(a_{4} = 9\)，求 \(x\) 的取值范围；
\item 设 \(\{a_{n}\}\) 是公比为 \(q\) 的等比数列，\(S_{n} = a_{1} + a_{2} + \dots + a_{n}\). 若 \(\frac{1}{3} S_{n} \leqslant S_{n+1} \leqslant 3S_{n}\)，\(n \in \mathbf{N}^{*}\)，求 \(q\) 的取值范围；
\item 若 \(a_{1}\)，\(a_{2}\)，\(\cdots\)，\(a_{k}\) 成等差数列，且 \(a_{1} + a_{2} + \cdots + a_{k} = 1000\)，求正整数 \(k\) 的最大值，以及 \(k\) 取最大值时相应数列 \(a_{1}\)，\(a_{2}\)，\(\cdots\)，\(a_{k}\) 的公差.
\end{enumerate}
\end{problem}
